% Body of the 250-word abstract (resumen_250_palabras.tex), centred on the
% package; mirrored in resumen_250_template.md. Numbers from numbers.tex.
Los mapas de evapotranspiración real ($\ETa$) obtenidos con modelos satelitales
de balance de energía, como METRIC, se usan para contabilizar y fiscalizar el
uso consuntivo del agua, pero sus elecciones de entrada rara vez se declaran:
lo habitual es forzar toda la escena con una sola estación meteorológica. Se
presenta EarthFlux, un paquete de R de código abierto que implementa METRIC con
física independiente del sensor, forzamiento meteorológico intercambiable
(estación representativa, regresión global, distancia inversa o regresión
geográficamente ponderada, GWR), validación cruzada espacial y diagnósticos que
marcan los valores anómalos sin corregirlos en silencio. En su versión
\EarthFluxVersion{}, \VerVerified{} de \VerRows{} formulaciones están
verificadas contra la fuente primaria y \NExpectations{} comprobaciones
automáticas respaldan la implementación. Como demostración, se aplicó a una
escena MODIS de la provincia de Cauquenes (Chile, 15-02-2019) con
\NStations{} estaciones públicas. Con las mismas anclas y ecuaciones, la sola
elección del forzamiento cambió la $\ETa$ media entre \DEtaCThreeGWRPct{} y
\DEtaCTwoIDWPct\,\%; con GWR, \VolDiffAbsCThreeGWR{} millones de m³ menos en un
día; el efecto, \RatioSA{} veces la incertidumbre de calibración, entró por la
$\ETr$ diaria. GWR no fue el mejor regionalizador: la
distancia inversa tuvo menor error fuera de muestra en \NCVDays{} días
(\CVTaCTwoIDWRmse{} frente a \CVTaCThreeGWRRmse~°C). El paquete es
experimental: faltan el flujo integrado y la descarga de múltiples escenas, y
la exactitud de $\ETa$ no se validó por falta de mediciones públicas de flujos.
Su aporte es hacer visible y auditable una incertidumbre de primer orden para
las decisiones sobre el agua.

\medskip\noindent\textbf{Palabras clave:} software científico abierto; evapotranspiración real;
METRIC; MODIS; regresión geográficamente ponderada; incertidumbre; gestión del agua.
